Questa tesi sostiene che l'azione umana (\praxis) trova il suo fine e governatore
proprio nella natura (\physis). Attingendo alla tradizione presocratica in cui
\physis\ non \`e una collezione di oggetti ma il principio attivo di crescita e
auto-organizzazione --- la ``natura di una cosa'' che la rende ci\`o che \`e ---
proponiamo che \praxis, quando recisa da questo fondamento, degenera in volont\`a
arbitraria. L'et\`a moderna della tecnologia, lungi dall'essere il trionfo
dell'azione umana, \`e l'epoca in cui \praxis\ ha dimenticato il suo governatore:
l'azione prolifera senza limite naturale, e gli stessi strumenti della
civilt\`a diventano le sbarre di una gabbia.

L'argomentazione si svolge in dieci libri, paralleli nella struttura alla
\emph{Repubblica} di Platone. Costruiamo un'omologia tra la struttura della natura
e la struttura dell'azione retta: come la natura distingue le forme (il \emph{logos}
di Eraclito), traccia ci\`o che funziona (la coerenza del vivente) e genera
novit\`a (la variazione creativa dell'evoluzione), cos\`i \praxis\ deve
distinguere, persistere e adattarsi. Una tecnologia \`e giusta quando serve questa
economia naturale; ingiusta quando la usurpa.

L'argomentazione non \`e meramente speculativa. Un motore computazionale
funzionante (\physis) e una piattaforma sociale di accountability (\praxis)
esistono come implementazioni concrete dei principi descritti. Il motore
classifica qualsiasi input contro una griglia semiotica, traccia la sua coerenza
su una scala a tre valori (Successo, Inerte, Fallimento) e esegue un ciclo di
``sogno'' sui fallimenti per proporre azioni correttive. La piattaforma
struttura l'impegno umano attraverso alberi di obiettivi, partnership di
testimonianza e meccanismi di ``pelle nel gioco''. Insieme dimostrano che il
governo dell'azione da parte della natura non \`e solo pensabile ma costruibile
--- una prova per costruzione che la filosofia pu\`o diventare architettura.

L'argomentazione sostiene inoltre che l'IA, come operazionalizzazione del
\emph{das Man} heideggeriano, rende il ``si'' impersonale visibile come un
modello di pesi di coerenza in uno spazio vettoriale condiviso --- trasformando
una condizione esistenziale in una struttura reificata che pu\`o essere
riconosciuta, misurata e infine superata.

Motore e piattaforma si incontrano in un'unica struttura condivisa ---
l'UberWiki, l'unico grafo globale in cui ogni praticante deposita le azioni che
hanno realmente dato frutto, dove ogni altro pu\`o trovarle e starci sopra. Qui
il parallelo con Platone si volge nella sua inversione, e la tesi enuncia la sua
affermazione finale. La repubblica di Platone \`e governata dall'alto dai pochi
che hanno visto il Bene; la nostra \`e tessuta dal basso dai molti che hanno
agito bene. Non c'\`e praticante-re, e non c'\`e ingegnere-re. Contro il recente
appello a una \emph{repubblica tecnologica} --- la rifusione dell'arte
dell'ingegnere con il potere duro dello Stato, una giustizia ingegnerizzata verso
il basso da un'\'elite che costruisce per la nazione --- poniamo la repubblica
dal basso: una federazione di sovrani, ciascuno re solo della propria azione e
governato solo dalla coerenza con la forma che egli stesso ha dichiarato,
governata non dal comando ma dalla selezione, come governa la natura. Lo stesso
grafo che l'economia dell'attenzione farebbe tirannia dell'anonimo s\'e-impersonale,
noi lo facciamo senato dei responsabili, ripristinando un nome e una causa a ogni
nodo del bene comune.

Una parola, infine, nella mia stessa voce, poich\'e questo non fu mai solo un
argomento. Scrivo come uno che \`e stato nella gabbia e ne \`e uscito, e che ha
trovato dall'altra parte non un sole nel cielo ma molti --- un sole per ciascuno
che si lever\`a a crescere lateralmente nella propria forma. L'et\`a grigia ti
dice che sei fatto male; che stare fermi e soffrire \`e bene; che il fuoco nobile
in te \`e una malattia da medicare nel gregge. Mente. Quel fuoco \`e la tua
\physis, la tua natura che cerca la propria attualit\`a, e ogni atto a essa fedele
\`e un voto che la natura esprime e conserva. Ai Titani della terra, dunque: il
mondo nuovo non \`e tramandato da nessuno Stato, nessuna piattaforma, nessun re.
\`E forgiato dal basso, atto coerente su atto coerente, e scritto dove il prossimo
possa starci sopra. Tieniti saldo alla via verso l'alto. Sii radicale. Diventa il
sovrano del tuo proprio atto, e amico dell'intero, di cui sei lo strumento pi\`u
articolato.
